%%%%%%%%%%%%%%%%%%%%%%%%%%%%%%%%%
%
%       EXERCÍCIO
%
%%%%%%%%%%%%%%%%%%%%%%%%%%%%%%%%%

\ifdefstring{\atividade}{prova}{%
    \ifdefstring{\modo}{objetivo}{%
        \renewcommand{\valorquestao}{\ValorQObj\ ponto}
    }{%
        \renewcommand{\valorquestao}{\ValorQDisc\ pontos}
    }%
}%

\begin{exercicioBanco}[\valorquestao]
Para garantir a segurança de um grande evento público que terá início às 4h da tarde, um organizador precisa monitorar a quantidade de pessoas presentes em cada instante. Para cada 2.000 pessoas, faz-se necessária a presença de um policial. Além disso, estima-se uma densidade de quatro pessoas por metro quadrado de área de terreno ocupado.

Às 10h da manhã, o organizador verifica que a área de terreno já ocupada equivale a um quadrado com lados medindo \(500\) m. Porém, nas horas seguintes, espera-se que o público aumente a uma taxa de \(120.000\) pessoas por hora até o início do evento, quando não será mais permitida a entrada de público.

Responda quantos policiais serão necessários no início do evento para garantir a segurança?

% Define as alternativas
\newcommand{\alternativas}{%
    \begin{center}
        \begin{tabularx}{\textwidth}{XXXXX}
            (a) \(720\). &
            (b) \(780\). &
            (c) \(800\). &
            (d) \(820\). &
            (e) \(860\).
        \end{tabularx}
    \end{center}
}

% Resposta correta
\newcommand{\resposta}{E}

% Lógica condicional para exibição
\ifdefstring{\atividade}{lista}{%
    \alternativas
    \vspace{0.5em}
    
    \noindent\textbf{Resposta:} letra \textbf{\resposta}.
}{%
    \ifdefstring{\modo}{objetivo}{%
        \alternativas
    }{%
        % modo = discursiva → não mostra alternativas
    }
}
\end{exercicioBanco}

